% СПО ЛКС - проектирование пользовательских протоколов
% Пынькин Д.А. (с) 2011
\documentclass[ignorenonframetext, hyperref={pdftex, unicode}]{beamer}

\usepackage[russian]{babel}
\usepackage[utf8]{inputenc}
\usepackage[T1]{fontenc}

\usepackage{beamerthemesplit}

%\usetheme{Pittsburgh}
%\usecolortheme{dolphin}

\usetheme{Copenhagen}
\usecolortheme{wolverine}


\usepackage{ulem}
%\usepackage{html}

\usepackage{verbatim}

\usepackage{ucs}
\usepackage{listings}
\lstloadlanguages{C, make, bash}

\lstset{escapechar=`,
	extendedchars=false,
	language=C, 
	tabsize=2, 
	columns=fullflexible, 
%	basicstyle=\scriptsize,
	keywordstyle=\color{blue}, 
	commentstyle=\itshape\color{brown},
%	identifierstyle=\ttfamily, 
	stringstyle=\mdseries\color{green}, 
	showstringspaces=false, 
	numbers=left, 
	numberstyle=\tiny, 
	breaklines=true, 
	inputencoding=utf8x,
	keepspaces=true,
	morekeywords={u\_short, u\_char, u\_long, in\_addr}
	}

\definecolor{darkgreen}{cmyk}{0.7, 0, 1, 0.5}

\lstdefinelanguage{diff}
{
    morekeywords={+, -},
    sensitive=false,
    morecomment=[l]{//},
    morecomment=[s]{/*}{*/},
    morecomment=[l][\color{darkgreen}]{+},
    morecomment=[l][\color{red}]{-},
    morestring=[b]",
}

\title[СПОЛКС (http://goo.gl/32cTB)]{Системное программное обеспечение\\локальных компьютерных сетей}
\author{Денис Пынькин}
\date{2011 -- 2012}
%\institution[БГУИР]{Белорусский государственный университет информатики и радиоэлектроники}
%\logo{\includegraphics[width=1cm]{logo-kafEVM.png}}


\subtitle{Форматы данных \\и проектирование пользовательских протоколов}

\mode<all>{\begin{frame}
\titlepage
\begin{center}
e-mail: denis.pynkin@bsuir.by\\
\end{center}
\begin{center}
{\bfseries http://goo.gl/32cTB}

{\tiny СЧАСТЬЕ ДЛЯ ВСЕХ, ДАРОМ, И ПУСТЬ НИКТО НЕ УЙДЕТ ОБИЖЕННЫЙ!\\
(c)Стругацкие, Пикник на обочине}
\end{center}
\end{frame}
}

%
% Далее начинается сама презентация
%

\section{Представление данных}

\begin{frame}{Представление данных}
	\begin{center}
	{\Large\bfseries Ясные протоколы лежат в основе хорошей практики.}
	\end{center}
\end{frame}

\begin{frame}{Схожесть задач}
	\begin{itemize}
		\item Проектирование форматов файлов для сохранения данных приложений в постоянном хранилище данных.
		\item Проектирование протоколов прикладного уровня для передачи (в том числе и через сеть) данных и команд между взаимодействующими программами.
	\end{itemize}
\end{frame}

\begin{frame}{Важное при проектировании прикладных протоколов}
	\begin{itemize}
		\item способность к взаимодействию; 
			\pause
		\item прозрачность;
			\pause
		\item расширяемость;
			\pause
		\item экономичность транзакций.
	\end{itemize}
\end{frame}

\begin{frame}{Типы данных}
	С машинной точки зрения удобно хранить данные сложной структуры в бинарном виде -- все поля имеют характерный для конкретной машины формат и указатели в виде адресов.

	\pause
	{\itshape Однако такие формы представления слабо подходят для длительного хранения и передачи.}
\end{frame}

\begin{frame}{Сериализация данных}
	Часто ЯП и/или платформа предоставляет средства для сериализации данных и восстановления исходной структуры из байтового потока.
\end{frame}

\begin{frame}{Минусы сериализации данных при использовании в протоколах}
	\begin{itemize}
		\item проблемы взаимодействия между машинами с разной архитектурой;
			\pause
		\item непрозрачность для других средств;
			\pause
		\item для сетевых протоколов часто целесообразно представлять структуру данных не в виде одного большого двоичного объекта,  а в виде последовательности транзакций или сообщений,  которые могут быть отклонены принимающей машиной.
	\end{itemize}
\end{frame}


\begin{frame}{Когда выгодно использовать бинарное представление?}
	\begin{itemize}
		\item Необходимость передачи больших блоков данных и разработчик формата действительно позаботился о достижении максимальной плотности полезной информации в потоке.

			Пример: мультимедиа форматы.
			\pause
		\item Существует жесткое ограничение времени и/или инструкций,  необходимых для интерпретации данных.

			Пример: протоколы сетевого уровня.
	\end{itemize}
\end{frame}

\begin{frame}{Текстовое представление данных}
	Дуг Макилрой (изобретатель pipes):
	\begin{itemize}
		\item Будьте готовы к тому,  что вывод каждой программы станет вводом другой,  еще неизвестной программы. 
			\pause
		\item Не загромождайте вывод посторонней информацией. 
			\pause
		\item Избегайте строгих табличных или двоичных форматов ввода. 
			\pause
		\item Не настаивайте на интерактивном вводе.
	\end{itemize}
\end{frame}

\begin{frame}{Текстовые потоки -- универсальный формат передачи данных}
	Они просты для чтения,  записи и редактирования человеком без использования специальных инструментов.
	\pause

	{\itshape Текстовые форматы -- прозрачны.}
\end{frame}

\begin{frame}{Текстовые потоки -- универсальный формат передачи данных}
	Если необходима производительность,  то можно внедрить сжатие текстового потока на уровне выше или ниже.

	Часто такая конструкция является более производительной,  чем бинарное представление.
\end{frame}

\begin{frame}{Текстовые потоки -- универсальный формат передачи данных}
	Упрощается интерпретация и анализ взаимодействия приложений,  а также написание тестовых программ.
\end{frame}

\begin{frame}{Текстовые потоки -- универсальный формат передачи данных}
	Серверные процессы часто запускаются с помощью суперсерверов подобных inetd/xinetd,  так что сервер получает команды на стандартный ввод и отправляет ответ на стандартный вывод.
\end{frame}

\begin{frame}{Текстовые потоки -- универсальный формат передачи данных}
	Можно взаимодействовать с сервером или клиентом с помощью программ {\tt telnet} или {\tt netcat}.
\end{frame}

\begin{frame}{Пример: SMTP}
\scriptsize
	{\bfseries telnet smtp.mail.ru 25}\\
	\color{gray}
	Trying 94.100.177.1...\\
	Connected to smtp.mail.ru.\\
	Escape character is '\^]'.\\
	\color{black}
	S: 220 smtp3.mail.ru ESMTP ready\\
	{\bfseries C: HELO my.server.org}\\
	S: 250 smtp3.mail.ru\\
	{\bfseries C: MAIL FROM: test\_testerson@mail.ru}\\
	S: 250 2.0.0 OK\\
	{\bfseries C: RCPT TO: denis.pynkin@bsuir.by}\\
	S: 250 Accepted\\
	{\bfseries C: DATA}\\
	S: 354 Enter message,  ending with "." on a line by itself\\
	{\bfseries C: from: test\_testerson@mail.ru\\
	C: to: denis.pynkin@bsuir.by\\
	C: subject: test\\
	C: Test message\\
	C: .}\\
	S: 250 OK id=1PmNuu-0006QX-00\\
	{\bfseries C: QUIT}
\normalsize
\end{frame}



\begin{frame}[fragile]{Пример: POP3}
\scriptsize
	{\bfseries telnet pop.mail.ru 110}\\
	S: +OK\\
	{\bfseries C: USER test\_testerson@mail.ru}\\
	S: +OK\\
	{\bfseries C: PASS test01}\\
	S: +OK Welcome!\\
	{\bfseries C: STAT}\\
	S: +OK 2 95232\\
	{\bfseries C: LIST}\\
	S: +OK 2 messages (95232 octets)\\
	S: 1 40269\\
	S: 2 54853\\
	S: .\\
	{\bfseries C: RETR 1}\\
	S: +OK 40269 octets\\
	{\color{gray}  Here is the message}\\
	{\bfseries C: DELE 1}\\
	S: +OK message 1 deleted\\
	{\bfseries C: QUIT}\\
	S: +OK POP3 server at mail.ru signing off\\
\normalsize
\end{frame}


\begin{frame}{Пример: FTP}
Управляющее соединение FTP:\\
\medskip
\tiny
	{\bfseries telnet ftp.mgts.by 21}\\
	S: 220 Welcome to ByFly FTP service.\\
	{\bfseries C: USER anonymous}\\
	S: 331 Please specify the password.\\
	{\bfseries C: PASS test\_testerson@mail.ru}\\
	S: 230 Login successful.\\
	{\bfseries C: PASV}\\
	S: 227 Entering Passive Mode (86, 57, 151, 3, 47, 76)\\
	\% ip=86.57.151.3 port=47*256+76=12108\\
	{\bfseries C: LIST}\\
	\%\%\%\%\% Здесь устанавливается информационное соединение №1 \%\%\%\%\%\\
	S: 150 Here comes the directory listing.\\
	S: 226 Directory send OK.\\
	{\bfseries C: PASV}\\
	S: 227 Entering Passive Mode (86, 57, 151, 3, 191, 44)\\
	{\bfseries C: RETR README}\\
	\%\%\%\%\% Здесь устанавливается информационное соединение №2 \%\%\%\%\%\\
	S: 150 Opening BINARY mode data connection for README (197 bytes).\\
	S: 226 File send OK.\\
	{\bfseries C: QUIT}\\
	S: 221 Goodbye.\\

\normalsize
\end{frame}

\begin{frame}{Пример: FTP}
Информационное соединение №1:\\
\medskip
\scriptsize
{\bfseries telnet ftp.mgts.by 12108}\\
\begin{verbatim}
-rw-r----r----    1 0    0       0 Apr 19  2010 2ban\_me.html\\
-rw-r----r----    1 0    0       197 Jan 27  2010 README\\
lrwxrwxrwx    1 0    0       18 Jan 27  2010 backports.org -> pub/backports.org/\\
drwxr-xr-x    2 0    0       23 May 20  2010 blog\\
drwxr-xr-x   10 0    0       4096 Feb 07 08:48 byfly\\
\end{verbatim}
\normalsize
\end{frame}

\begin{frame}{Пример: FTP}
Информационное соединение №2:\\
\medskip
\scriptsize
{\bfseries telnet ftp.mgts.by 48940}\\
\begin{verbatim}
Welcome to ByFly public archive (ftp.byfly.by)\\
	located at MGTS,  Minsk,  Belarus\\
	sponsored by Beltelecom (www.beltelecom.by)\\
	2 x Intel(R) Xeon(R) CPU X3210  @ 2.13GHz\\
	8 Gb RAM,  6 TB SATA Storage\\
\end{verbatim}
\normalsize
\end{frame}

\begin{frame}{}
\Huge
\begin{center}
	Спасибо за внимание!
	\bigskip
	Вопросы?
\end{center}
\normalsize
\end{frame}

\mode<all>{\begin{frame}{}
\Huge
\begin{center}
	Спасибо за внимание!
	\bigskip
	Вопросы?
\end{center}
\normalsize
\end{frame}
}

\end{document}
%Конец файла
